%=====================================================================
\chapter{Discounts, Vouchers and Referrals}\label{ch:discounts}
%=====================================================================
\locator{3}{Part III --- From Basket to Order}

\begin{outcomes}
By the end of this chapter you will be able to:
\begin{tight}
  \item \textbf{Distinguish} the four price-reduction devices a storefront
        offers, and \textbf{choose} the right one for a stated objective.
  \item \textbf{Explain} why order totals are an ordered pipeline, and
        \textbf{predict} what moving a line in it does to the customer's bill.
  \item \textbf{Compute} the margin consequence of a discount, and
        \textbf{state} the extra volume needed to pay for it.
  \item \textbf{Describe} how a referral is tracked and paid, and
        \textbf{identify} what can be claimed twice.
  \item \textbf{Report} honestly on a feature a prescribed outline names and a
        platform does not provide.
\end{tight}
\end{outcomes}

\begin{prereq}
\chref{shipping}: tax classes, rates and rules, and the worked order total.
\chref{checkout}: the checkout sequence, and that an order copies its totals.
A working checkout, so that discounts can be observed on real orders.
\end{prereq}

\begin{keyterms}
discount $\bullet$ quantity break $\bullet$ special price $\bullet$ coupon
$\bullet$ coupon code $\bullet$ gift voucher $\bullet$ store credit $\bullet$
reward points $\bullet$ affiliate $\bullet$ referral $\bullet$ commission
$\bullet$ tracking code $\bullet$ totals pipeline $\bullet$ sort order
$\bullet$ gross margin
\end{keyterms}

\section{Four Devices, and What Each One Is For}

``Give the customer money off'' is one sentence and four different mechanisms,
which differ in \emph{who qualifies} and \emph{when the shop finds out}.

\begin{definitionbox}[Quantity break, Special, Coupon, Credit]
A \term{quantity break} reduces the unit price above a threshold --- buy ten,
pay less each. It is attached to the product, applies to everyone, and the
customer learns of it by filling the basket.

\smallskip
A \term{special} is a dated price that replaces the usual one for a period, for
a customer group. Also attached to the product; also automatic; visible on the
product page as a struck-through old price.

\smallskip
A \term{coupon} is a code the customer must type. It is attached to nothing
--- it is its own object, with its own dates, its own usage limits and its own
restriction to particular products or categories. The code is the point: it
lets the shop give a discount \emph{to the people it gave the code to}.

\smallskip
\term{Store credit} is a balance on a customer's account that is spent like
money. It is not a discount at all; it is a liability the shop already owes.
\end{definitionbox}

\begin{conceptbox}[Which to reach for]
\begin{tight}
  \item \textbf{To move stock} --- a special. It needs no action from the
        customer and no code to leak.
  \item \textbf{To reward larger baskets} --- a quantity break, which pays for
        itself in picking and delivery.
  \item \textbf{To measure a campaign} --- a coupon, and \emph{a different code
        per channel}. The code is the only part of a discount that tells you
        where the customer came from.
  \item \textbf{To settle a complaint} --- store credit, which costs the shop
        its margin rather than its price list, and brings the customer back.
\end{tight}

\smallskip
The mistake is to use a coupon for everything. A public code is a price cut
with extra steps, and it teaches customers to search for a code before buying
--- which is a cost that arrives later and never goes away.
\end{conceptbox}

\section{The Clause the Platform Does Not Implement}

HEC's outline names ``Discounts, \emph{Vouchers}, and Referrals''. A gift
voucher --- a code bought by one person as a stored-value instrument and spent
by another --- is a distinct thing from a coupon, and the shop in this course
does not have one.

\begin{alertbox}[There are no gift vouchers in this platform]
Checked directly rather than assumed:

\begin{tight}
  \item \textbf{No tables.} \texttt{SHOW TABLES LIKE '\%voucher\%'} returns
        nothing. The database has 159 tables and not one of them is a voucher.
  \item \textbf{No totals extension.} The nine shipped are \texttt{sub\_total},
        \texttt{shipping}, \texttt{coupon}, \texttt{tax}, \texttt{credit},
        \texttt{reward}, \texttt{handling}, \texttt{low\_order\_fee} and
        \texttt{total}. No voucher.
  \item \textbf{One trace remains:} a demonstration image called
        \texttt{gift-voucher-birthday.jpg}. The feature existed in version 3
        and was removed; the picture was not.
\end{tight}

\smallskip
\textbf{Why this is in the handout rather than quietly skipped.} HEC's clause
list is the contents page of a book published in 2010 (see the Course
Specification), when stored-value vouchers were a standard cart feature. The
outline has not been revised; the software has. Stating that plainly is more
useful to you than either pretending the feature exists or pretending the
clause does not.

\smallskip
\textbf{What to do if a client asks for vouchers.} Store credit
(\texttt{oc\_customer\_transaction} and the \texttt{credit} total) already
provides the spending half: a balance that draws down across orders. What is
missing is the \emph{buying} half --- selling a voucher as a product, issuing a
code, and letting a different person redeem it. That is a module, and
\chref{modules} is where you would write it. It is the single most realistic
piece of commissioned work in this course.
\end{alertbox}

\section{Totals Are a Pipeline, and Position Decides the Bill}

\chref{shipping} computed an order total by hand and found that the VAT line
included the shipping charge. That was not arithmetic; it was \emph{ordering}.
Every total line is produced by a separate extension with a \term{sort order},
and each one works on the running total left by the lines before it.

\dsfig{%
\begin{tikzpicture}[
  ln/.style={rounded corners=2pt, draw=#1, line width=1.0pt, fill=#1!8,
             text=#1!60!black, font=\bfseries\scriptsize, align=center,
             text width=2.0cm, minimum height=0.72cm},
  num/.style={font=\scriptsize\bfseries, text=neutral},
  arr/.style={-{Stealth[length=1.9mm]}, primary!50, line width=0.8pt}
]
  \node[ln=secondary] (a) at (0,0)    {Sub-Total};
  \node[ln=neutral]   (b) at (2.35,0) {Reward};
  \node[ln=neutral]   (c) at (4.70,0) {Shipping};
  \node[ln=greenhl]   (d) at (7.05,0) {Coupon};
  \node[ln=accent]    (e) at (9.40,0) {Tax};
  \node[ln=secondary] (f) at (11.75,0) {Total};

  \foreach \a/\b in {a/b, b/c, c/d, d/e, e/f}{\draw[arr] (\a) -- (\b);}
  \foreach \n/\x in {1/0, 2/2.35, 3/4.70, 4/7.05, 5/9.40, 9/11.75}{
    \node[num] at (\x,0.62) {\n};}

  \node[font=\scriptsize, text=greenhl, align=center, text width=3.0cm]
       at (7.05,-1.15) {The discount is subtracted \emph{here}\dots};
  \node[font=\scriptsize, text=accent, align=center, text width=3.0cm]
       at (9.40,-1.15) {\dots so tax is charged on what is left.};
  \node[font=\scriptsize\itshape, text=neutral, align=center, text width=11.5cm]
       at (5.9,-2.1) {Swap the two numbers and every page still works, every
       line still adds up, and the customer pays more.};
\end{tikzpicture}}%
{The totals pipeline as the shop ships it, with each extension's sort order
above it. The numbers are editable fields in the admin panel, and nothing
warns you what they mean.}{totalspipeline}

\begin{worked}
\textbf{The same basket, with the coupon on either side of the tax.} Two real
orders, produced by \texttt{code/m12\_discounts.py --sortorder}. The basket is
\chref{checkout}'s: one MacBook at \$500.00 and two iPhones at \$101.00. The
coupon is the one the shop ships with, code \texttt{2222}, 10\% off.

\smallskip
\textbf{As shipped --- coupon at 4, tax at 5:}

\smallskip
\begin{tabular}{@{}L{6.4cm}L{3.4cm}L{4.2cm}@{}}
Sub-Total & 702.00 & \\
Flat Shipping Rate & 5.00 & \\
Coupon (2222) & $-70.20$ & $10\% \times 702.00$ \\
Eco Tax (-2.00) & 6.00 & $3 \times 2.00$ \\
VAT (20\%) & 127.36 & $20\% \times 636.80$ \\
\textbf{Total} & \textbf{770.16} & \\
\end{tabular}

\smallskip
The VAT base is $702.00 + 5.00 - 70.20 = 636.80$: the discount has already
happened, so tax is charged on what the customer actually pays.

\smallskip
\textbf{Coupon moved to 6, after tax:}

\smallskip
\begin{tabular}{@{}L{6.4cm}L{3.4cm}L{4.2cm}@{}}
Sub-Total & 702.00 & \\
Flat Shipping Rate & 5.00 & \\
Eco Tax (-2.00) & 6.00 & \\
VAT (20\%) & 141.40 & $20\% \times 707.00$ \\
Coupon (2222) & $-70.20$ & \\
\textbf{Total} & \textbf{784.20} & \\
\end{tabular}

\smallskip
\textbf{The difference is \$14.04}, and it is exactly 20\% of the \$70.20
discount --- tax charged on money the customer was never asked for. On a \$770
order that is 1.8\%, produced by changing one integer in an admin field that
offers no explanation of what it does.

\smallskip
\textbf{Both orders are internally consistent.} Both add up. Both display
correctly. No validation objects, because there is nothing malformed about
either --- one of them is simply wrong about the world.
\end{worked}

\begin{alertbox}[A sequencing rule the checkout does not tell you about]
Applying a coupon \emph{after} the customer has chosen a payment method
invalidates the choice. The coupon changes the order total, the stored payment
method no longer matches the recomputed list, and the final step answers:

\smallskip
\hspace*{1em}\texttt{\{"error":"Payment method is incorrect!"\}}

\smallskip
--- after which no order is created. The code was accepted, every earlier step
reported success, and the sale is lost at the last click.

\smallskip
This was found by driving the sequence from a script, exactly as in
\chref{checkout}, and it is the second fault in this course that is invisible
from the browser. The remedy in your own shop is to put the coupon field
\emph{before} the method steps, or to re-select the method whenever a total
changes.
\end{alertbox}

\section{What a Discount Actually Costs}

A discount is not a marketing decision with a financial side effect. It is a
financial decision, and the arithmetic is short enough to do before the
meeting rather than after it.

\begin{worked}
A product sells for \textbf{Rs~1{,}000} and costs the shop \textbf{Rs~700}.

\smallskip
\textbf{Gross margin today:} $1{,}000 - 700 = \mathbf{Rs~300}$ per unit, 30\%.

\smallskip
\textbf{With a 10\% coupon:} the customer pays $Rs~900$; the cost is unchanged.
\[ 900 - 700 = \mathbf{Rs~200} \]

\smallskip
\textbf{The discount was 10\% of the price and 33\% of the profit.}
\[ \frac{300 - 200}{300} = 33.3\% \]

\smallskip
\textbf{How much extra volume pays for it?} To earn the same total profit:
\[ \frac{300}{200} = 1.5 \]
\textbf{The shop needs 50\% more orders} to stand still. Not 10\% more --- 50\%.

\smallskip
And that is the optimistic case, because it assumes every discounted sale is a
sale that would not otherwise have happened. In practice some share of the
customers who used the code would have bought anyway, and every one of those is
pure loss. If half of them would have bought anyway, the required uplift roughly
doubles again.

\smallskip
\textbf{The general rule:} at a gross margin of $m$ and a discount of $d$ (both
as fractions of price), the volume multiple needed to break even is
\[ \frac{m}{m - d} \]
which goes to infinity as $d$ approaches $m$, and is \emph{negative} --- that
is, unachievable at any volume --- beyond it. A 30\% discount on a 30\% margin
cannot be made up by selling more.
\end{worked}

\section{Referrals: Paying Someone Else to Send Customers}

HEC's third term. The platform's mechanism is the \term{affiliate}: a customer
account flagged as one, given a \term{tracking code}, and paid a commission on
orders that arrive carrying it.

\begin{examplebox}[How a referral is tracked, in four steps]
\begin{tightnum}
  \item A customer applies at \adminpath{Account, Affiliate} and is approved
        --- or not, depending on \uiel{Affiliate Approval} in the settings.
  \item They receive a tracking code and link their audience to
        \texttt{?tracking=\textit{code}}.
  \item A visitor arriving on that link has the code stored. The default
        commission in this shop is \textbf{5\%}, set at
        \adminpath{System, Settings, Edit, Option}.
  \item When an order completes, a commission transaction is written against
        the affiliate's account, becoming store credit or a payable.
\end{tightnum}

\smallskip
\textbf{The question that decides whether this works:} how long does the
tracking code survive? If it lasts one request, the affiliate is paid only for
customers who buy immediately. If it lasts a month, the affiliate is paid for
customers who were going to buy anyway. Neither extreme is right, and the shop
--- not the affiliate --- must choose.
\end{examplebox}

\begin{alertbox}[Where referral money leaks]
\begin{tight}
  \item \textbf{Self-referral.} An affiliate buys through their own link. A 5\%
        commission is a 5\% personal discount, paid for by the shop and
        recorded as marketing spend.
  \item \textbf{Stacking with a coupon.} A 10\% code plus a 5\% commission is
        15\% off a margin that may be 30\%, and half the profit is gone with
        two independently reasonable decisions.
  \item \textbf{Commission on refunded orders.} Unless the commission is
        reversed when the order is, the shop pays for a sale that was undone
        --- \chref{orders}'s subject, and the commonest unreconciled figure in
        a small shop's accounts.
  \item \textbf{Commission on tax and shipping.} Paying 5\% of \$854.40 rather
        than 5\% of the \$702.00 of goods means paying the affiliate a share of
        money that belongs to the tax authority and the courier.
\end{tight}
\end{alertbox}

\begin{pitfall}
\begin{tight}
  \item \textbf{Treating totals sort order as cosmetic.} It moved a bill by
        \$14.04 and nothing warned.
  \item \textbf{Putting the coupon field after the payment step}, which loses
        the sale at the last click with an error nobody sees.
  \item \textbf{Assuming a demonstration coupon is usable.} The shipped codes
        are disabled and expired --- which is correct of them, and the first
        thing to check when a code ``does not work''.
  \item \textbf{Approving a discount without the margin arithmetic.} 10\% off a
        30\% margin needs 50\% more orders.
  \item \textbf{Publishing one coupon code everywhere}, which makes the
        campaign unmeasurable and the discount permanent.
  \item \textbf{Paying commission on gross order value} rather than on goods.
  \item \textbf{Promising a client gift vouchers} because the outline mentions
        them. This platform has none; it is a module, not a setting.
\end{tight}
\end{pitfall}

\begin{lab}[Move one number, and change what the customer pays]
\begin{lstlisting}[language=bash]
# --- 1. Find out why the shipped coupon does not work ---------------
cd code
docker compose exec db mariadb -ushop -pshoppw opencart -e \
  "SELECT code,name,status,date_start,date_end FROM oc_coupon;"
# Two reasons, not one. Name both before fixing either.

# --- 2. Enable it ----------------------------------------------------
python m12_discounts.py --coupon on

# --- 3. Place one order with it and take the total apart ------------
python m12_discounts.py --order
# Reproduce every line with a calculator.

# --- 4. Predict, in writing --------------------------------------
# If the coupon were applied AFTER tax instead of before, would the
# customer pay more, less, or the same? By how much? Commit first.

# --- 5. Run the experiment ------------------------------------------
python m12_discounts.py --sortorder

# --- 6. Account for the difference exactly --------------------------
# It is not an arbitrary number. Show where it comes from.

# --- 7. Build a measurable campaign ---------------------------------
# Admin: Marketing > Coupons. Make THREE codes for the same 10% offer,
# one per channel. Say how you would later tell which worked.

# --- 8. Do the margin arithmetic before you would ship it -----------
# Your product costs 700 and sells for 1000. Work out the volume
# uplift your 10% code needs. Then do it for 20%.

# --- 9. Put everything back -----------------------------------------
python m12_discounts.py --restore
\end{lstlisting}

\smallskip
\textbf{What to notice.} Step 4 against step 5. Most people predict ``the
same'' --- the discount and the tax are both there either way. The 1.8\%
difference is the whole chapter, and it lives in a numeric field with no label
explaining it.
\end{lab}

\begin{checkpoint}
\begin{tightnum}
  \item The same basket produced totals of 770.16 and 784.20 depending on one
        setting. Explain what the setting was, why the difference is exactly
        14.04, and which total is correct --- giving your reason in terms of
        what tax is owed on.
  \item A client with a 25\% gross margin wants to run a permanent 15\% ``loyal
        customer'' code. Work out the volume uplift needed to break even, state
        two reasons the real requirement is higher than your figure, and
        propose an alternative that costs less.
  \item HEC's outline requires vouchers and the platform has none. Describe
        what you would tell a client who asked for them, what the platform
        already provides, and what specifically you would have to build.
\end{tightnum}
\end{checkpoint}

\begin{chaptersummary}
\begin{tight}
  \item Four devices, differing in who qualifies and when the shop learns of
        it: quantity breaks and specials are automatic and attached to
        products; a coupon is a code, which is what makes a campaign
        measurable; store credit is a liability, not a discount.
  \item This platform has \textbf{no gift vouchers} --- no tables, no totals
        extension, only a leftover demonstration image. HEC's clause reflects a
        2010 book, not current software. Say so rather than fudging it; the
        gap is a module, and \chref{modules} is where it would be written.
  \item Order totals are an ordered pipeline. Each line works on what the
        previous lines left, and every extension's position is an editable
        integer with no explanation attached.
  \item Worked: the same basket totalled 770.16 with the coupon before tax and
        784.20 with it after --- a difference of 14.04, exactly 20\% of the
        70.20 discount. Both are internally consistent; one is wrong about what
        tax is owed on.
  \item Applying a coupon after the payment method is chosen invalidates the
        choice and loses the order at the last click, with an error the browser
        does not show.
  \item Worked: a 10\% discount on a 30\% margin removes 33\% of the profit and
        needs 50\% more orders to break even --- $m/(m-d)$ in general, and
        unachievable once $d$ reaches $m$.
  \item A referral is tracked by a code and paid as commission. The leaks are
        self-referral, stacking with coupons, commission on refunded orders,
        and commission on tax and shipping.
\end{tight}
\end{chaptersummary}

\begin{reviewq}
  \item \textbf{(MCQ)} Moving the coupon line after the tax line in the totals
        pipeline causes the customer to pay: (a)~less; (b)~the same;
        (c)~more, by the tax on the discount; (d)~more, by the discount.
  \item \textbf{(MCQ)} A 10\% discount on a product with a 30\% gross margin
        removes what share of the profit? (a)~10\%; (b)~about a third;
        (c)~half; (d)~it depends on the volume.
  \item \textbf{(Short)} Give two things a coupon can do that a special cannot,
        and one thing a special does better.
  \item \textbf{(Short)} Explain why a shop should issue a different coupon
        code for each channel, even for the same offer.
  \item \textbf{(Short)} Explain why paying affiliate commission on the order
        grand total rather than on the goods value is a mistake, using the
        worked order in \chref{shipping} for your figures.
  \item \textbf{(Applied)} A shop runs a 20\% code and a 5\% affiliate scheme
        on products with a 35\% gross margin. Work out the margin on an order
        that uses both, say what the shop should do about it, and state how you
        would enforce your answer in the platform.
  \item \textbf{(Applied)} Design the gift voucher the platform lacks. State
        the tables you would add, how a voucher is bought, how it is redeemed
        by a different person, what happens when it part-pays an order, where
        it sits in the totals pipeline and why, and the two ways it could be
        spent twice.
\end{reviewq}
